\section{The Three Trials}

\begin{quotex}
To accept suffering in itself and for itself, to consent to it, to seek it, to love it, to make its identifying mark and the very object of disinterested and generous love, to put perfect action in sorrowful passion, to be active up to the point of death, to make of every act a death and of death itself the act par excellence, here is the triumph of the will that disconcerts again nature and that in fact engenders in man a new --- and more than human --- life.
\flright{\textsc{Maurice Blondel}\footnote{\url{http://www.iep.utm.edu/blondel/\#H3}}, \emph{Action}}
\end{quotex}


As a young man, Julius Evola embarked on an intense course of autodidactic studies of German Idealism, French Personalism, and Oriental Metaphysics, particularly Taoism and Tantrism. Around the same time he discovered Rene Guenon and read what was available from him up until that point. The result was his philosophy of \textbf{Magical Idealism}, which is succinctly summarized in the \emph{The Individual and the Becoming of the World}. The argument starts from the self-awareness as an acting person and is developed logically and thoroughly until its conclusion. However, one cannot be convinced by the argument on a strictly intellectual level, on which it is neither true nor false. Only by going through a process of self-transformation can one determine the validity or falseness of magical idealism. In his book, \emph{Saggi sull'Idealismo Magico} (``Essays on Magical Idealism'') EMI , published in 1925, Evola explains:

\begin{quotex}
The truth or falsity of idealism---and that means if man can or can't give certainty and sense to his life and his experience---cannot be demonstrated theoretically: it can be decided not through an intellectual act, but through a concrete realization.
\end{quotex}


In EMI, Evola not only provides a logical development of magical idealism, but he also proposes a methodology whereby the reader can achieve such a concrete realization. There are two phases, the first being the state of purification or catharsis. The first stage requires some preparation which involves \textbf{three trials: fire, suffering and love}. The goal of this preparation is:

\begin{enumerate}


\item To confirm the absolute autonomy of the center of the I

\item To generate the principle of an entirely active acting
\end{enumerate}

\paragraph{Trial by Fire}


The trial by fire begins with the experience of negation. The I habitually pulls its support from peripheral elements such as science, culture, attachments, faith, and so on. In this trial, he must destroy every such mental and emotional support. He must ``deny every faith, violate every moral and social law, scorn every sentiment of humanity, every love and generosity, every passion, affirm an implacable and all-pervasive skepticism, reaching finally a conscious and critical madness.''

The horror and disvalue that is instinctively felt by such precepts, proceeds only from an internal unconscious fear, that one is not sufficient in himself. Evola gets support for this view from the Tantric teaching of bhutashuddi\footnote{\url{http://www.itimes.com/public/people/iti133802/blog/Dhyan-Yoga-by-Chakra-Meditation}} or ``purification from the elements'' and in alchemy to the ``liberation of the metals''.

\paragraph{Trial by Suffering}


The trial by suffering is the condition of remaining close to the denial of one's own life, but insofar as the denial is not authored by the I, which is freed from the dependence of the object to deny. In the trial by fire, the individual made himself independent from those determinations in a relative way: he needs them to deny them, hence is dependent on them. In this trial, he is freed only by detaching from himself, eliminating the negative power: not willing it, not attributing it to himself, but rather simply undergoing it, accepting it as something foreign and transcending his own will.

From this perspective, we see the value in Stoicism and Christian resignation; this is how we understand why various saints invoke suffering as a divine grace.

\paragraph{Trial by Love}


This trial is no longer about the abstract denial of the self, but in that deeper denial of the self, that is the existence in oneself of a thing as object of unconditioned love. At this point it is not about destruction, but of constructing oneself at every moment, through a renewed act of love and renunciation, on a plane higher than oneself where there is the impassibility of the Watcher\footnote{\url{http://www.kktanhp.com/vipassana\_htm.htm}} or, better, the Lord, outside any storm or tumult whatsoever, whether interior or exterior. It is not a question of the denial of passion, but rather an indifference that has no need to exclude anything. This is the greatest non-resistance alluded to by Lao Tzu.

In the other trials, the I was autonomous as pure essence; in this trial, as pure act. There is an action that has its origins in the desire for something, not in itself, but only as it relates to the I. This reveals an insufficient centre. But to want something for itself, is to remove the I as object of volition, and thus to really have it.

Equally, violent and passionate action against things revels they have a reality for the I, which is the opposite of absolute self-determination. Violating things is tantamount to violating only oneself. The fundamental principle of magic is:

\begin{quotex}
To really want a thing, it is necessary to want it not for the I but for itself, that is, to love it.
\end{quotex}


Violence is the way of weakness and impotence; love and gentleness, the way of strength and mastery. This is the profound Taoist doctrine: not to desire to have in order to have, to give in order to possess, to cede in order to dominate, to sacrifice oneself in order to realize oneself. This is wei-wu-wei, or ``acting without acting''.

From this we see the profound meaning and value of the greatest teachings on humility, submission, self-denial, detachment, surrendering one's own will to God. The elimination of every pride, a life permeated with humility and self-denial, the constant death of one's own will are tasks to which the individual is inadequate unless he manifests a power infinitely superior to that demand for any madness of negation and destruction. This is a hard trial that only constitutes the preliminary condition for the life of the ruler.

\paragraph{Commentary}


The preceding text on the three trials are taken from EMI; I don't always use quotes because the material may have been abridged or rearranged.

This text should convince everyone by now that we are not engaged in an abstract intellectual discussion, but rather in a personal, existential process of self-transformation. This should also show the Traditional sources of Evola's system: Tantra, Hermetism, Gnosticism, Catholicism, Stoicism, Taoism. These two points should no longer be an issue.

It is important to remember that these three trials constitute just the preliminary part of the first phase of purification; more will be forthcoming in future posts.

The Trial by Fire may sound harsh until one realizes that nothing real is being denied. Instead, the I is freeing itself from artificial thought constructions that he had previously considered to be real. This is consistent with Traditional practices. In Tomberg's system, it is part of the purification of consciousness described in Arcanum II.

The Trial by Suffering is the acceptance of whatever appears to be. The individual's attitude towards suffering must be no different from his attitude toward a rainy day; it is something exterior and of no concern. Besides the writings of the saints and Stoics that Evola mentions, there is also Nietzsche's concept of \emph{amor fati}.

With the Trial by Love, Evola makes clear that he is not speaking of egoic desires. Desirelessness is the end result of this Trial. This involves a state of inner detachment; things are willed not to fulfill my needs (which would be a limitation) but rather for their own sake. Drawing on the Tao Te Ching, Evola compares this actionless action to Taoist ideas. He also likens it to Arcanum XI of the Tarot, ``Strength'', showing his familiarity with Hermetism. It is also constructive to compare Evola's depiction of the Ruler who rules without violence with what Tomberg writes about the Emperor in Meditations on the Tarot.

This trial is creative; while the Trial by Fire was eliminative. This Trial is active, willing things from Love. Evola points out that love, respect, devotion and so on are not ends in themselves but rather the means to a higher power.

\begin{quotex}
Morality is a means, never an end: it is not of value in itself, but in as much as through it, the will can empower its own affirmation.
\end{quotex}



\flrightit{Posted on 2011-03-12 by Cologero}

\begin{center}* * *\end{center}

\begin{footnotesize}
\begin{sffamily}
\texttt{Albion on 2011-03-12 at 21:55 said:}

Very clarifying, vis Evola. ``To become God'' is to imply that one is NOT God now, therefore, only to be taken literally (in becoming) as it is stated and meant. Thank you. The metaphor of the ``rainy day'' is a very good one.

\hfill

\texttt{Cologero on 2011-03-13 at 11:58 said:}

Evola begins this chapter with this enigmatic quote from John 10:34

``Is it not written in your law: I said you are gods?''

So he is attempting to bring out the possible meaning for that.

Evola also claims that is is just incorporating elements from Yoga (Tantra/Patanjali), the various Western esoteric traditions (including Christian), and Theosophy (Blavatsky, Steiner). He claims to be purifying them (e.g., ignoring any claims of special paranormal powers as in Steiner), reducing them to their inner meaning, that is, to what can actually be experienced in consciousness. So it is by this standard that he is to be judged: has he, in fact, synthesized a method for spiritual development stripped of the particularities of their historical and cultural accretions? Or has he somehow distorted them?

In one of the early reviews of the book, a Theosophist (Jasink) points out that from Oriental conceptions, it makes sense, in understanding the I as Atman and Brahman. That is how Evola needs to be read, at least for anyone familiar with the terminology and teachings of the Vedanta. However, he goes on to say that it becomes ``muddied or objectionable'' was translated into Western languages. That is why Atman and Brahman are always left untranslated, a practice which really obscures their meaning. Jasink explains that in the West, I has the meaning of the empirical ego, and this I is certainly not the creator of everything, as there are multiple subjectivities in this sense.

It should be clear that Evola is not referring to this empirical Ego, nor is the Will the mere satisfaction of the desire of that Ego to dominate. Absolutely to the contrary, such an egoic life demonstrates a weakness, not power.

\hfill

\texttt{Albion on 2011-03-25 at 07:59 said:}

Does Evola attempt to ``order'' these trials, or are they occuring ``as they occur''?

\hfill

\texttt{Albion on 2011-03-25 at 08:09 said:}

Cologero, do you think you could comment (at some point in time) on Christian Platonism and Plato in general? If where you are headed is true north (and I think it is) then situating it in relation to Plato would be helpful to the rest of us (I think), as it is a familiar landmark. Then again, I am having to rethink everything I thought I knew about Plato.

\hfill

\texttt{Cologero on 2011-03-25 at 12:29 said:}

I assume by ``as they occur'' you mean ``as the opportunity arises.'' But there are always opportunities.

Sample opportunities:

1) Do you question the sources of your thoughts and opinions. Do your ideas prevent you from hearing anything ``new''.

2) How do you react to negative people or events in life? Do you automatically react, or is there a ``space'' between the stimulus and response where free will can act.

3) How do weakness and concupiscence affect your life? Are you frustrated when the satisfaction of desires is delayed? Do fears inhibit your acts?

There are plenty of opportunities in the course of the day. These are areas that I think are best explored in the private area of the forums.

\hfill

\texttt{Cologero on 2011-03-25 at 12:39 said:}

The place to debate specific points about Plato or some other specific writer is in the forums.

However, I am planning to bring in more about the cubist painter Albert Gleizes, who was quite familiar with Guenon and employed Traditional themes in his work. Interestingly, he considered Plotinus, Augustine and Boethius (all Platonists) as the intellectual fathers of the Medieval era.

Evola considered Plotinus (along with Patanjali) as secure guides to the inner exploration of the Spirit, what he termed ``metaphysical positivism''. These are connections that need to be explored.

\hfill

\texttt{Albion on 2011-03-27 at 01:28 said:}

Thank you, that is exactly what I was asking. I will see you in the Forum.

\hfill

\texttt{Paulo on 2014-08-23 at 22:31 said:}

``Violence is the way of weakness and impotence; love and gentleness, the way of strength and mastery. This is the profound Taoist doctrine: not to desire to have in order to have, to give in order to possess, to cede in order to dominate, to sacrifice oneself in order to realize oneself. This is wei-wu-wei, or ``acting without acting''.

``Morality is a means, never an end: it is not of value in itself, but in as much as through it, the will can empower its own affirmation.''

The fear of pain,has to be understood,if the will is strong,love gonna appear as a anwser to pain and fear!

\hfill

\texttt{aegishjalmer on 2014-10-09 at 13:07 said:}

The lust of the eyes, the lust of the flesh, and the pride of life

\hfill

\texttt{Frater M on 2017-03-13 at 22:27 said:}

If things are not willed for my own ``needs', once resolved via detachment, but still willed ``for their own sake'', one might then ask, why will them at all then?

What is this other ``own sake'' that obligates further willing whatsoever? Is it not still, all things considered, another limitation, another willing?

Only Bodhisattvic logic works here.

\hfill

\texttt{Noct on 2022-04-17 at 01:30 said:}

I do not think that wanting things for themselves implies a heteronomy or something that ``obligate'', rather the opposite: it is not an imposition, but a free act of the will. Remaining in the fullness of the Self, one experiences a kind of ``superabundance of being'' and desires in a positive sense, not as something that is suffered, but as a creative activity, the highest good of all beings.

Is this not the teaching of the Lord when he says that we are to love our neighbor as ourselves?

\hfill

\texttt{Cologero on 2022-04-17 at 09:01 said:}

Noct, I do not think that we are reading the same text, since you draw the opposite conclusion. ``It is about constructing oneself at every moment, through a renewed act of love and renunciation, on a plane higher than oneself.'' That sounds creative to me.

You can love your neighbor as yourself, but do you do so to show what a great guy you are (i.e., to enhance the ``I''), or because you actually love the neighbor.

\hfill

\end{sffamily}
\end{footnotesize}

